\qquad Este trabalho tem como objetivo implementar e analisar cinco
algoritmos clássicos de ordenação por comparação estudados na
disciplina de Projeto de Algoritmos: Insertion Sort, Selection Sort,
Bubble Sort, Shell Sort e Merge Sort. Foi desenvolvida uma aplicação
em C++, com criação automática das pastas e arquivos pelo próprio
programa, capaz de gerar instâncias de teste de diferentes tamanhos
(10, 100, 1.000, 10.000, 100.000 e 1.000.000 de elementos) e
diferentes características (crescente, decrescente e randômica),
executar cada algoritmo sobre elas e registrar o tempo de execução
observado.

\qquad Os resultados foram organizados em arquivos de entrada, saída e
tempo para cada algoritmo, consolidados em tabelas e gráficos
individuais e comparativos que permitem analisar o comportamento de
cada algoritmo isoladamente e compará-los entre si. As medições
confirmam o comportamento teórico esperado para cada algoritmo:
desempenho linear no melhor caso do Insertion e do Bubble Sort
(entrada já ordenada), desempenho quadrático em qualquer caso para o
Selection Sort, e uma clara superioridade dos algoritmos de menor
complexidade: para 1.000.000 de elementos randômicos, o Merge Sort
(0,078 s) e o Shell Sort (0,106 s) foram cerca de mil a doze mil vezes
mais rápidos que os três algoritmos $O(n^2)$ (80,6 s a 938,7 s). O
Merge Sort, cujo custo é $\Theta(n \log n)$ independentemente da ordem
inicial dos dados, foi o mais rápido para entradas randômicas, ao passo
que o Shell Sort foi o mais rápido para entradas já ordenadas ou
decrescentes, evidenciando a forte influência da estratégia de cada
algoritmo --- e, no caso do Insertion, Selection e Bubble Sort, também
do grau de ordenação inicial dos dados --- sobre o desempenho observado.
